\documentclass{ohm_project_description}

\projecttitle{Fallstudie zum Einsatz von neuen Technologien in Unternehmen}
\projectauthor{TTZ Nürnberger Land / TH Nürnberg}
\projectdate{2026}

\begin{document}

\maketitle
\thispagestyle{fancy}

Im Rahmen der Fallstudie sollen primär kleinere und mittlere Unternehmen (KMU) im Nürnberger Land hinsichtlich der Anwendung neuer Technologien wie Künstliche Intelligenz, mobiler und stationärer Robotersysteme sowie der Möglichkeit der Automatisierung von Tätigkeiten und Prozessen befragt werden.

Dazu gehört die Ermittlung der bereits genutzten Technologien, der Bereitschaft neue Technologien in den Bereichen Produktion, Qualitätsmanagement und technische Kommunikation einzusetzen sowie der dabei auftretenden Herausforderungen und Hürden. Es sollen Umfragen methodisch vorbereitet, geeignete Unternehmen identifiziert und Entscheidungsträger in Interviews befragt und ausgewertet werden. Das Ergebnis ist eine zusammenfassende Darstellung der gewonnenen Erkenntnisse und eine kurze Einordnung hinsichtlich relevanter Technologien.

Das \textbf{Technologietransferzentrum Nürnberger Land} hat das Ziel, in der Region den Wissenstransfer aus der Forschung in überwiegend kleinere und mittlere Unternehmen voranzutreiben und verfügt über ein breit gefächertes Netzwerk in der Region.

\section*{Arbeitspakete}
\begin{itemize}[leftmargin=0.5cm]
    \setlength\itemsep{.1em}
    \item Methodische Vorbereitung und Durchführung von Unternehmensumfragen
    \item Identifikation geeigneter Unternehmen im Nürnberger Land
    \item Durchführung von Interviews mit Entscheidungsträgern
    \item Auswertung und zusammenfassende Darstellung der Ergebnisse
    \item Einordnung relevanter Technologien (KI, Robotik, Automatisierung)
\end{itemize}

\section*{Voraussetzungen}
\begin{itemize}[leftmargin=0.5cm]
    \setlength\itemsep{.1em}
    \item Interesse an Fallstudien und empirischer Forschung
    \item Breit gefächerter technischer Hintergrund
    \item Kontaktfreude und Organisationsfähigkeit
    \item Eigenständiges und zielgerichtetes Arbeiten
\end{itemize}

\vspace{0.5cm}
Das Thema ist als \textbf{Projektarbeit} konzipiert.

\vfill
\textcolor{ohm_red}{\rule{\linewidth}{0.4mm}}
\textbf{\textcolor{ohm_red}{TTZ Nürnberger Land / Labor für mobile Robotik}} \\
\begin{tabular}{@{}ll}
\textbf{Betreuer:} & M.\,Sc. Patrick Fußy \\
\textbf{E-Mail:}   & \href{mailto:patrick.fussy@th-nuernberg.de}{patrick.fussy@th-nuernberg.de} \\
\end{tabular}

\end{document}
